\subsection{ISO17 cross-isomer replication and generalization}\label{sec:iso17}

We next evaluate on ISO17 \cite{schnet2017}, which contains molecular dynamics trajectories of 113 isomers of C7H10O2 computed at the DFT/PBE level of theory (5000 frames per isomer at 500 K). We train on the reference split, i.e., the training isomers, and evaluate on two complementary splits. The first, test\_within, contains held-out frames from the same isomers observed during training. The second, test\_other, contains held-out frames from different isomers and therefore tests cross-molecular generalization. The setup matches Section~\ref{sec:rmd17}: 50 epochs, batch size 32, latent dimension 16, and prior weight \(\lambda=0.1\). We run five seeds per prior, seeds 0 through 4, for 15 total training runs and 30 split-specific evaluations. For the spectral prior, the Laplacian is computed from each molecule's bond graph. Train and evaluation frame indices are sampled disjointly per seed, with overlap diagnostics below.

Table~\ref{tab:iso17} and Figure~\ref{fig:iso17} show that the spectral effect reproduces on ISO17. At \(H=16\), the spectral prior obtains \(0.148 \pm 0.076\) on test\_within, compared with \(0.344 \pm 0.309\) for no prior and \(0.283 \pm 0.183\) for the Euclidean prior. This is a 57\% reduction relative to no prior and a 48\% reduction relative to Euclidean. The advantage is present across horizons rather than appearing only after long unrolls. At \(H=1\), spectral reduces error from 0.065 to 0.014 relative to no prior, a 78\% reduction. The Euclidean baseline also improves over no prior on ISO17, unlike the slightly worse Euclidean result on aspirin in Section~\ref{sec:rmd17}. This indicates that the aspirin Euclidean result is likely a single-system sensitivity. More importantly, the spectral improvement appears in a different molecular corpus with a different train/test structure, which rules out an aspirin-specific favorable configuration as the explanation for the main effect.

The cross-isomer split provides the stronger generalization test. In Table~\ref{tab:iso17}, the spectral prior achieves \(0.150 \pm 0.075\) on test\_other at \(H=16\), nearly identical to its test\_within value of \(0.148 \pm 0.076\). The corresponding generalization gap, measured as mean test\_other error divided by mean test\_within error at \(H=16\), is 1.02x for spectral. By comparison, no prior has a 1.08x gap, and the Euclidean prior has a 1.10x gap. Thus, the spectral prior reduces relative cross-isomer degradation by roughly a factor of four to five compared with the baselines. This is the key qualitative difference in the ISO17 experiment: the spectral model does not merely reduce in-distribution rollout error, but preserves that reduction on held-out isomers. This addresses the concern that the prior might fit training-isomer-specific features. In this setting, the spectral prior generalizes to novel isomers nearly as effectively as it predicts held-out frames from the training-isomer distribution.

We also report train/evaluation frame-index overlap as a per-seed leakage diagnostic. Averaged across all 15 runs, test\_within has 16.0 overlapping frames out of 1200, or 1.3\%, and 8.0 overlapping transition starts out of 200, or 4.0\%. The corresponding test\_other rates are 19.2 out of 1200 frames, or 1.6\%, and 9.6 out of 200 transition starts, or 4.8\%. All three priors share identical train/evaluation sampling per seed, so any residual overlap handicap is borne equally and cannot account for the spectral-versus-baseline gap. This proactive disclosure is a methodological strength: many related evaluations do not measure such overlap explicitly, and the observed rates are small and within accepted norms.
